\subsection{向量空间与线性系统}
我们现在重新考察线性方程组与高斯消元法，
所借助的是本章的工具与术语。
我们要说明三点。

先说第一点，
回想第一章中的一个洞见：
高斯消元法的工作方式是对行取
线性组合\Dash 若两个矩阵
通过行变换
$A\longrightarrow\cdots\longrightarrow B$ 相联系，则
$B$ 的每一行都是 $A$ 的各行的线性组合。
因此，研究一般的行变换、
特别是高斯消元法的合适框架，是下述向量空间。

\begin{definition} \label{df:RowSpace}
%<*df:RowSpace>
矩阵的\definend{行空间\/}\index{matrix!row space}\index{row space}
是其行集合的张成。
\definend{行秩\/}\index{row!rank}\index{matrix!row rank}\index{row rank}
是这个空间的维数，即线性无关的行的数目。
%</df:RowSpace>
\end{definition}

\begin{example}
若
\begin{equation*}
  A=\begin{mat}[r]
          2  &3  \\
          4  &6
       \end{mat}
\end{equation*}
则 \( \rowspace{A} \) 是二分量行向量空间中
的这个子空间。
\begin{equation*}
   \set{c_1\cdot\rowvec{2 &3}+c_2\cdot\rowvec{4  &6}
          \suchthat c_1,c_2\in\Re }
\end{equation*}
第二个行向量线性依赖于第一个，
因此可以把上面的描述化简为
$\set{c\cdot\rowvec{2 &3}\suchthat c\in\Re }$。
\end{example}

\begin{lemma}     \label{le:RowSpUnchByGR}
%<*lm:RowSpUnchByGR>
若两个矩阵 \( A \) 与 \( B \) 通过一次行变换
\begin{equation*}
  A\grstep{\rho_i\leftrightarrow\rho_j}B
  \quad\text{或}\quad
  A\grstep{k\rho_i}B
  \quad\text{或}\quad
  A\grstep{k\rho_i+\rho_j}B
\end{equation*}
（其中 $i\neq j$ 且 $k\neq 0$）相联系，则它们的行空间相等。
从而，行等价的矩阵有相同的行空间，因此
也有相同的行秩。
%</lm:RowSpUnchByGR>
\end{lemma}

\begin{proof}
%<*pf:RowSpUnchByGR0>
推论 One.III.\ref{cor:RowsOfEqMatsLinCombos}
表明，当 \mbox{$A\longrightarrow B$} 时，$B$ 的每一行
都是 \( A \) 的各行的线性组合。
也就是说，按上面的术语，\( B \) 的每一行
都是 $A$ 的行空间中的元素。
于是 $\rowspace{B}\subseteq\rowspace{A}$ 成立，因为集合
$\rowspace{B}$ 的成员是 $B$ 的各行的线性组合，从而它是
$A$ 的各行的组合的组合，
再由线性组合引理，它也是 $\rowspace{A}$ 的成员。
%</pf:RowSpUnchByGR0>

%<*pf:RowSpUnchByGR1>
对于另一个方向的包含关系，
回想引理 One.III.\ref{le:RowOpsRev}：行变换是可逆的，
所以
\mbox{$A\longrightarrow B$} 当且仅当
\mbox{$B\longrightarrow A$}。
于是与上一段同样地可得
$\rowspace{A}\subseteq\rowspace{B}$。
%</pf:RowSpUnchByGR1>
% Thus the two sets are equal.
\end{proof}

当然，高斯消元法系统地执行这些行变换，
其目标就是阶梯形。

\begin{lemma}  \label{le:RowsEchMatLI}
%<*lm:RowsEchMatLI>
阶梯形矩阵的非零行
构成一个线性无关的集合。
%</lm:RowsEchMatLI>
\end{lemma}

\begin{proof}
%<*pf:RowsEchMatLI>
引理~One.III.\ref{le:EchFormNoLinCombo}
说的是，阶梯形矩阵的任一非零行都不是
其余各行的线性组合。
这个结果用本章的术语重述了该结论。
%</pf:RowsEchMatLI>
\end{proof}

于是，用本章的语言来说，
高斯消去法的工作方式是消除各行之间的线性相关关系，
而保持张成不变，直到非零行之间
不再剩下非平凡的线性关系。
简言之，高斯消元法给出行空间的一个基。

\begin{example}
对任一矩阵，我们可以这样得到行空间的一个基：
执行高斯消元法，然后取所得的阶梯形
矩阵的非零行。
例如，
\begin{equation*}
  \begin{mat}[r]
    1  &3  &1  \\
    1  &4  &1  \\
    2  &0  &5
  \end{mat}
  \grstep[-2\rho_1+\rho_3]{-\rho_1+\rho_2}
  \repeatedgrstep{6\rho_2+\rho_3}
  \begin{mat}[r]
    1  &3  &1  \\
    0  &1  &0  \\
    0  &0  &3
  \end{mat}
\end{equation*}
就给出行空间的基 $\sequence{\rowvec{1 &3 &1},
            \rowvec{0 &1 &0},
            \rowvec{0 &0 &3} }$。
它既是起始矩阵的也是终止矩阵的行空间的基，
因为这两个行空间相等。
\end{example}

利用这一技巧，我们还能对那些不直接涉及行向量的
张成求出基。

\begin{definition} \label{df:ColumnSpace}
%<*df:ColumnSpace>
矩阵的\definend{列空间}\index{column!space}\index{matrix!column space}
是其列集合的张成。
\definend{列秩}\index{column!rank}\index{rank!column}
是列空间的维数，即线性无关
列的数目。
%</df:ColumnSpace>
\end{definition}

我们对列空间的兴趣源于对线性方程组的研究。
一个例子是，方程组
\begin{equation*}
  \begin{linsys}{3}
    c_1  &+  &3c_2  &+  &7c_3  &=  &d_1  \\
   2c_1  &+  &3c_2  &+  &8c_3  &=  &d_2  \\
         &   &c_2   &+  &2c_3  &=  &d_3  \\
   4c_1  &   &      &+  &4c_3  &=  &d_4
  \end{linsys}
\end{equation*}
有解当且仅当 \( d \) 组成的向量是其余各列向量的
线性组合，
\begin{equation*}
  c_1\colvec[r]{1 \\ 2 \\ 0 \\ 4}
  +c_2\colvec[r]{3 \\ 3 \\ 1 \\ 0}
  +c_3\colvec[r]{7 \\ 8 \\ 2 \\ 4}
  =\colvec{d_1 \\ d_2 \\ d_3 \\ d_4}
\end{equation*}
也就是说，\( d \) 组成的向量位于系数矩阵的
列空间中。

\begin{example}   \label{ex:BasisForColSpace}
给定这个矩阵，
\begin{equation*}
  \begin{mat}[r]
    1  &3  &7  \\
    2  &3  &8  \\
    0  &1  &2  \\
    4  &0  &4
  \end{mat}
\end{equation*}
为求列空间的一个基，
先暂时把列变成行并作化简。
\begin{equation*}
    \begin{mat}[r]
       1  &2  &0  &4  \\
       3  &3  &1  &0  \\
       7  &8  &2  &4
    \end{mat}
  \grstep[-7\rho_1+\rho_3]{-3\rho_1+\rho_2}
  \repeatedgrstep{-2\rho_2+\rho_3}
  \begin{mat}[r]
     1  &2  &0  &4  \\
     0  &-3 &1  &-12\\
     0  &0  &0  &0
  \end{mat}
\end{equation*}
现在再把行变回列。
\begin{equation*}
  \sequence{
     \colvec[r]{1 \\ 2 \\ 0 \\ 4},
     \colvec[r]{0 \\ -3 \\ 1 \\ -12} }
\end{equation*}
结果就是所给矩阵的列空间的一个基。
\end{example}

\begin{definition} \label{df:Transpose}
%<*df:Transpose>
矩阵的\definend{转置}\index{transpose}\index{matrix!transpose}
是把它的行与列互换
所得的结果，使得
矩阵 \( A \) 的第~\( j \) 列是 \( \trans{A} \) 的第~\( j \)
行，反之亦然。
%</df:Transpose>
\end{definition}
\noindent 于是我们可以把前面的例题总结为「转置，
化简，再转置回来」。

我们甚至可以——以容忍向量空间「相同」
这一尚且模糊的想法为代价——
用高斯消元法来求其他类型的向量空间中张成的基。

\begin{example}
为求空间 \( \polyspace_4 \) 中
\( \set{x^2+x^4,2x^2+3x^4,-x^2-3x^4} \) 的张成的一个基，
把这三个多项式
看成与行向量 \( \rowvec{0 &0 &1 &0 &1} \)、
\( \rowvec{0 &0 &2 &0 &3} \) 及
\( \rowvec{0 &0 &-1 &0 &-3} \) 「相同」，并应用高斯消元法
\begin{equation*}
  \begin{mat}[r]
    0  &0  &1  &0  &1  \\
    0  &0  &2  &0  &3  \\
    0  &0  &-1 &0  &-3
  \end{mat}
  \grstep[\rho_1+\rho_3]{-2\rho_1+\rho_2}
  \repeatedgrstep{2\rho_2+\rho_3}
  \begin{mat}[r]
    0  &0  &1  &0  &1  \\
    0  &0  &0  &0  &1  \\
    0  &0  &0  &0  &0
  \end{mat}
\end{equation*}
然后再翻译回去，得到基
\( \sequence{x^2+x^4,x^4} \)。
（正如前面提到的，
我们将在下一章的开头把「相同」
这个说法讲清楚。）
\end{example}

这样，本小节的第一点就是
本章的工具让我们对高斯消去法
有了更具概念性的理解。

至于第二点，
注意到对矩阵作行变换可能会改变其列空间。
\begin{equation*}
  \begin{mat}[r]
    1  &2  \\
    2  &4
  \end{mat}
  \grstep{-2\rho_1+\rho_2}
  \begin{mat}[r]
    1  &2  \\
    0  &0
  \end{mat}
\end{equation*}
左边矩阵的列空间包含第二个分量
非零的向量，
而右边矩阵的列空间
只包含第二个分量为零的向量，所以这两个空间
是不同的。
这一观察使得接下来的结果出人意料。

\begin{lemma} \label{le:RowOpsNoChngColRnk}
%<*lm:RowOpsNoChngColRnk>
行变换不改变列秩。
%</lm:RowOpsNoChngColRnk>
\end{lemma}

\begin{proof}
%<*pf:RowOpsNoChngColRnk>
换个说法，若 $A$ 化简为 $B$，
则 $B$ 的列秩等于 $A$ 的列秩。

如果我们能证明行变换不影响各列之间的线性关系，
这个证明即告完成，因为列秩就是
最大的无关列集合的大小。
也就是说，我们将证明：各列之间存在某个关系
（例如第五列等于第二列的
两倍加上第四列）当且仅当在行变换之后
该关系仍然存在。
而这正是本书的第一个定理，
定理 One.I.\ref{th:GaussMethod}：~在各列间的
关系
\begin{equation*}
  c_1\cdot\colvec{a_{1,1} \\ a_{2,1} \\ \vdots \\ a_{m,1}}
   +\dots+
  c_n\cdot \colvec{a_{1,n} \\ a_{2,n} \\ \vdots \\ a_{m,n}}
   =\colvec[r]{0 \\ 0 \\ \vdotswithin{0} \\ 0}
\end{equation*}
中，行变换使解集 \( (c_1,\ldots,c_n) \) 保持不变。
%</pf:RowOpsNoChngColRnk>
\end{proof}

要说明高斯消元法不仅对行空间、
也对列空间能有所断言，
另一种方式是高斯–若尔当消元。
它终止于矩阵的
简化阶梯形，如下所示。
\begin{equation*}
  \begin{mat}[r]
    1  &3  &1  &6  \\
    2  &6  &3  &16 \\
    1  &3  &1  &6
  \end{mat}
  \grstep{}\cdots\grstep{}
  \begin{mat}[r]
    1  &3  &0  &2  \\
    0  &0  &1  &4  \\
    0  &0  &0  &0
  \end{mat}
\end{equation*}
考虑这个结果的行空间与列空间。

本小节前面讲的第一点告诉我们，
为求行空间的一个基，只需收集那些含有
首主元的行。
然而，由于这是简化阶梯形，
求列空间的一个基同样容易：~收集那些含有首主元的
列，
\( \sequence{\vec{e}_1,\vec{e}_2} \)。
% (Linear independence is obvious.
% As to span, the other columns are in the span
% since they all have a third component of zero.)
于是，对简化阶梯形矩阵，
我们可以用本质上相同的方法求出行空间与列空间的基，
即取矩阵中含有首主元的
那些部分——行或列。

\begin{theorem} \label{th:RowRankEqualsColumnRank}
%<*th:RowRankEqualsColumnRank>
对任意矩阵，
行秩与列秩相等。
%</th:RowRankEqualsColumnRank>
\end{theorem}

\begin{proof}
%<*pf:RowRankEqualsColumnRank0>
把矩阵化成简化阶梯形。
此时
行秩等于首主元的个数，因为它等于
非零行的数目。
同时，首主元的个数
也等于列秩，因为含有首主元的那些列组成的集合
是标准基中的一些 \( \vec{e}_i \)，
并且该集合线性无关，且张成
全体列所成的集合。
因此，在简化阶梯形矩阵中，行秩等于列
秩，因为二者都等于首主元的个数。
%</pf:RowRankEqualsColumnRank0>

%<*pf:RowRankEqualsColumnRank1>
但引理 \nearbylemma{le:RowSpUnchByGR} 与 \nearbylemma{le:RowOpsNoChngColRnk}
表明，用行变换化成简化阶梯形时，行秩与列秩
都不改变。
从而，原矩阵的行秩与列秩也相等。
%</pf:RowRankEqualsColumnRank1>
\end{proof}

\begin{definition} \label{df:Rank}
%<*df:Rank>
矩阵的\definend{秩\/}\index{rank}就是它的行秩或列秩。
%</df:Rank>
\end{definition}

于是，我们在本小节中提出的第二点就是：
矩阵的列空间与行
空间有相同的维数。

我们的最后一点是：在研究向量空间时
自然出现的那些概念，
恰恰就是我们在研究线性方程组时已考察过的那些。

\begin{theorem}  \label{th:RankVsSoltnSp}
%<*tm:RankVsSoltnSp>
对于 \( n \) 个未知数、系数矩阵为
\( A \) 的线性方程组，下述命题
\begin{tfae}
   \item \( A \) 的秩为 \( r \)
   \item 相应齐次方程组的解空间具有
     维数 \( n-r \)
\end{tfae}
是等价的。
%</tm:RankVsSoltnSp>
\end{theorem}

\par\noindent 于是，若方程组至少有一个特解，
则对于解
集而言，参数的个数等于 $n-r$，即变量个数
减去系数矩阵的秩。

\begin{proof}
%<*pf:RankVsSoltnSp>
\( A \) 的秩为 \( r \) 当且仅当对 \( A \) 作高斯消去法
最终得到 \( r \) 个非零行。
这当且仅当与 \( A \) 行等价的阶梯形矩阵
有 \( r \) 个首变量时成立。
而这又当且仅当有 \( n-r \) 个自由变量时成立。
%</pf:RankVsSoltnSp>
\end{proof}

\begin{corollary} \label{co:EquivToNonsingular}
%<*co:EquivToNonsingular>
设矩阵 $A$ 为 \( \nbyn{n} \)，下述命题
\begin{tfae}
  \item \( A \) 的秩为 \( n \)
  \item \( A \) 非奇异
  \item \( A \) 的各行构成线性无关的集合
  \item \( A \) 的各列构成线性无关的集合
  \item 系数矩阵为 \( A \) 的任何线性方程组都有且
    仅有一个解
\end{tfae}
是等价的。
%</co:EquivToNonsingular>
\end{corollary}

\begin{proof}
%<*pf:EquivToNonsingular>
显然 \( \text{(1)}\iff\text{(2)}\iff\text{(3)}\iff\text{(4)} \)。
最后一条，\( \text{(4)}\iff\text{(5)} \)，成立是因为 \( n \)
个列向量构成的集合线性无关当且仅当它是
\( \Re^n \) 的一个基，而方程组
\begin{equation*}
  c_1\colvec{a_{1,1} \\ a_{2,1} \\ \vdots \\ a_{m,1}}
   +\dots+
  c_n\colvec{a_{1,n} \\ a_{2,n} \\ \vdots \\ a_{m,n}}
   =\colvec{d_1 \\ d_2 \\ \vdots \\ d_m}
\end{equation*}
对所有 \( d_1,\dots,d_n\in\Re \) 的选取都有唯一解，当且仅当
左边由诸 \( a \) 构成的向量组成一个基。
%</pf:EquivToNonsingular>
\end{proof}

\begin{remark}\cite{Munkres}  \label{rem:MunkresCommment}
有时，人们会把本小节的结果
误记为：\( m \)~个方程、\( n \)~个未知数的方程组的通解
使用 \( n-m \) 个参数。
方程的个数并不是相关的数；真正要紧的
是独立方程的个数，即极大无关
集中方程的个数。
当有 \( r \) 个独立方程时，通解用到
\( n-r \) 个参数。
\end{remark}


